\section{Regular torsors and delooping}
\label{sec:torsors}

\subsection{Left torsors and the regular example}
For a set-level group $G$, a left torsor is a left $G$-action on a type $X$
which is merely inhabited and has contractible transporter types
\begin{equation}
  \sum_{g:G}(g\cdot x=y)
\end{equation}
for all $x,y:X$.  The second condition is the type-theoretic free-and-
transitive condition; the truncation in the first condition is important,
since a general torsor has no selected origin.
\formal{\texttt{src/TorsorDelooping/LeftTorsor.agda}}
{\texttt{IsLeftTorsor}; \texttt{LeftTorsor}}
{The fields are \agda{inhabited : ∥ X ∥₁} and the contractible-transporter field \texttt{regular}.}

The regular left torsor is $G$ itself under left multiplication.  It is
inhabited by $1$, and its transporter is contractible through the imported
right-translation equivalence.  The use of right translations is compatible
with the left action: it proves the needed regularity without changing the
chosen action convention.
\formal{\texttt{src/TorsorDelooping/Regular.agda}}
{\texttt{regularLeftAction}; \texttt{regularLeftTorsor}}
{Regularity uses \texttt{·GroupAutomorphismR}; the inhabitant is the group unit.}

\subsection{The univalent moduli type}
The generic construction first records a left action, then adds the
propositional fields asserting sethood, the action laws, mere inhabitation,
and contractible transporters.  The Structure Identity Principle therefore
identifies paths of structured torsors with equivariant equivalences.  It
constructs a moduli type $G\mathrm{Tors}_{L,\ell_X}$ and an equivalence
\begin{equation}
 \operatorname{TorsorEquiv}(T,U)\simeq(T=U).
 \label{eq:sip-torsors}
\end{equation}
\formal{\texttt{src/TorsorDelooping/Moduli.agda}}
{\texttt{LeftTorsorModuli}; \texttt{TorsorEquiv}; \texttt{torsorPathEquiv}; \texttt{fromLeftTorsor}}
{\texttt{torsorPathEquiv} is the SIP equivalence; \texttt{fromLeftTorsor} embeds the regular-torsor record into the moduli type.}

\subsection{The constructive complex basepoint}
Set $G=\CHunits$.  The local units module supplies the regular torsor
\agda{complexUnitsRegularLeftTorsor}; its image in the moduli type defines
\begin{align}
 \BG\CHunits&\coloneqq\CHunits\mathrm{Tors}_{L,0},\\
 b_{\CHunits}&\coloneqq\text{the regular left $\CHunits$-torsor}.
\end{align}
The source records \agda{BComplexUnits : Type₁}; that universe level is
inherited from the generic structured-type construction.
\formal{\texttt{src/ConstructiveComplex/Units.agda}}
{\texttt{complexUnitsRegularLeftTorsor}}
{This is the regular torsor of \texttt{ComplexUnitsGroup}.}
\formal{\texttt{src/ConstructiveComplex/Delooping.agda}}
{\texttt{BComplexUnits}; \texttt{bComplexUnits}}
{Both are direct specializations of \texttt{LeftTorsorModuli} and \texttt{torsorBase}.}
